\section{The cost of simultaneous torsion selection}

Let $v$ be the valuation on $\overline{\mathbb Q}_p$ with $v(p)=1$ and $v(0)=+\infty$. Write $\mathcal U_h$ for the roots of exact order $p^h$, $e_h=(p-1)p^{h-1}$, and $r_h=1/e_h$. For an integral polynomial $g=\sum_j a_jT^j$, its Gauss valuation at radius parameter $r$ is
\[
 G_r(g)=\min_j\{v(a_j)+jr\}.
\]
Thus $v(g(y))\ge G_r(g)$ for $v(y)=r$. The root of unity in the source selector is evaluated in the coordinate $T=\zeta-1$.

\paragraph{The question.}
The closing research question (ii) of the incoming \emph{Quantitative twists} manuscript \cite{priortorsion} asks whether a bound on the shell exponent $H$ can be polynomial in the total budget and inverse tolerance, uniformly in degrees and packet size, while controlling total positive excess. In the earlier manuscript on torsion selection the relevant one-variable assertion has the following parameters. Each polynomial $g_i$ has an integral algebraic coefficient set and a witness $x_i$ with $v(x_i)\ge\rho$. Set $t_i=v(g_i(x_i))<\infty$ and assume $\sum_i t_i\le B$. One seeks $h\le H$ with $r_h\le\rho$ and $\zeta\in\mathcal U_h$ such that
\[
 \sum_i \big(v(g_i(\zeta-1))-t_i\big)_+\le\epsilon,
 \qquad (u)_+=\max(u,0).
\]
The argument from the source gives a cardinality-independent sufficient bound exponential in $(B+1)/\epsilon$ for $0<\epsilon<1$, with fixed depth and dimension. The quantitative twists manuscript gives a polynomial bound for a fixed finite packet with an exponent that depends on the number of constraints, for comparison to the fixed-depth Gauss valuations. Its example that vanishes on every point of $H$ shells costs budget $H$. Controlling positive excess permits cheaper obstructions.

\paragraph{Theorem (exponential lower bound).}
Fix $p=2$, $\rho=1$ and $0<\epsilon<1$. For every integer $H\ge1$ there are $H$ monic integral polynomials, all defined over $\mathbb Q_2(\zeta_{2^H})$, with the same witness $x_i=0$ and total budget
\[
 B_H\le 2+4\epsilon\sum_{j=1}^{H-1}\frac1j .
\]
For every $1\le h\le H$ and every $\zeta\in\mathcal U_h$, at least one of the individual excesses exceeds $\epsilon$. In particular for $0<\epsilon\le1/8$ and $B=3$ no selector can guarantee success by shell
\[
 H=\left\lfloor\exp\!\left(\frac{1}{8\epsilon}\right)\right\rfloor.
\]
Together with the harmonic upper bound this gives $\exp(\Theta(1/\epsilon))$ for the worst-case first successful shell exponent at fixed $B=3$. The order of the root itself is $2^h$.

\paragraph{Proof.}
Fix a primitive $2^h$-th root $z_h$ in each shell. For $1\le k\le h$, choose the transversal $A_k=\{1,3,\ldots,2^k-1\}$ of odd residue classes modulo $2^k$ and put
\[
 g_{h,k}(T)=\prod_{a\in A_k}\big(T-(z_h^a-1)\big),
 \qquad t_{h,k}=2^{k-h}.
\]
All $2^{k-1}$ roots have valuation $r_h=2^{1-h}$. Multiplicativity of Gauss valuations gives
\[
 v(g_{h,k}(0))=G_1(g_{h,k})=G_{r_h}(g_{h,k})=t_{h,k}.
\]
The identity $v(z_h^u-1)=2^{v_2(u)}/2^{h-1}$ for $2^h\nmid u$ follows from the order of $z_h^u$. Indeed, the primitive $2^m$-th roots give conjugate differences from $1$, whose product up to sign is $\Phi_{2^m}(1)=2$.

For any odd $b$ and each $1\le j<k$, exactly $2^{k-j-1}$ members of $A_k$ satisfy $v_2(a-b)=j$. These contribute $t_{h,k}$ to $v(g_{h,k}(z_h^b-1))$ for each such $j$. The one remaining $a$ agrees with $b$ modulo $2^k$. Its factor has valuation at least $2^k/2^{h-1}=2t_{h,k}$, with $+\infty$ allowed when the factor vanishes. Consequently
\[
 v\big(g_{h,k}(z_h^b-1)\big)-t_{h,k}\ \ge\ k\,t_{h,k}.
\]
When $k<h$ equality is attained at $b=1+2^k$. When $k=h$ every evaluation vanishes. The amplification by $k$ comes from the full set of nested congruence conditions on $a-b$.

For shell $h$ choose the least $k=k(h)\in\{1,\ldots,h\}$ such that $k2^{k-h}>\epsilon$. This exists because the expression at $k=h$ is $h$. Use the polynomial $g_{h,k(h)}$ as the $h$-th constraint. The preceding inequality blocks the entire shell. To sum the costs, the terms with $k(h)=1$ cost at most $\sum_{h\ge1}2^{1-h}=2$. When $k(h)=k\ge2$, minimality gives
\[
 2^{k-h}\le\frac{2\epsilon}{k-1}.
\]
With $k$ fixed these terms are distinct powers of $2$ and their sum is at most twice the largest. Summing over $k=2,\ldots,H$ yields the asserted budget. Finally $\sum_{j=1}^{H-1}j^{-1}\le1+\log H$, so the displayed exponential choice of $H$ costs at most $2+4\epsilon+1/2\le3$.\hfill$\square$

\paragraph{A single polynomial for the moving baseline.}
Set $P=\prod_{h=1}^H g_{h,k(h)}$. Then $G_1(P)=B_H$. At each primitive point of shell $h$,
\[
 v(P(\zeta-1))-G_{r_h}(P)>\epsilon.
\]
Indeed every factor has nonnegative excess over its own $G_{r_h}$, and the shell's designated factor has excess greater than $\epsilon$. This shows exponential sharpness in inverse tolerance for the source's stronger one-polynomial comparison to the moving Gauss valuation. Replacing that baseline by $G_1(P)$ can absorb an individual loss through savings on the other factors. This distinction accounts for the polynomial fixed-packet result in the earlier report.

\paragraph{What bounded ramification buys.}
Let $K/\mathbb Q_p$ be finite with absolute ramification index $e$. Suppose all $g_i$ are integral polynomials over $K$. Among the admissible shells $r_h\le\rho$, at most $\lfloor eB\rfloor$ contain a root valuation of any $g_i$ when $\sum_i G_\rho(g_i)\le B$. On every other admissible shell, \emph{all} primitive points simultaneously satisfy $v(g_i(\zeta-1))=G_{r_h}(g_i)\le G_\rho(g_i)$. Thus $\lfloor eB\rfloor+1$ consecutive admissible shells guarantee a shell of zero positive excess. The witness-budget version follows from $G_\rho(g_i)\le t_i$.

To see this, let $m_{ih}$ count the roots of $g_i$ of valuation $r_h$ with multiplicity. Their multiset is stable under the absolute Galois group of $K$. When $m_{ih}>0$ their product belongs to $K^\times$ and $m_{ih}r_h\ge 1/e$. Factorization of $g_i$ into linear factors gives
\[
 \sum_{h:r_h\le\rho} m_{ih}r_h\le G_\rho(g_i)-G_0(g_i)\le G_\rho(g_i).
\]
In fact a root of valuation $w$ contributes $\min(w,\rho)-\min(w,0)$ to the difference, a nonnegative amount. This argument includes roots at zero by the convention $w=+\infty$. Summation proves the exceptional-shell count. On a shell devoid of roots of matching valuation, the ultrametric inequality gives equality with the Gauss valuation in each linear factor. The counterexamples above use fields of growing ramification and freely chosen algebraic integral packets. Any application to narrower arithmetic families must separately determine which coefficient fields and polynomial packets can actually occur.
